\section{Bootstrapping the module loader}
\label{module-loader}
The language features described up to this point let us build a program
called a \emph{module-loader} that is able to process user programs
organised in modules such as the one defined as follows:
\begin{listingbox}\begin{ildcode}
(module
  (doc "this module calculates 5!")
  (exports main)
  (imports
    (host (macroexpand <= * +))
    ("core/prelude.ild" (if))
    ("core/module-utils.ild" (fn)))
  (defs
    (!fn main () (fact 5))

    (!fn fact (x)
      (!if (<= x 0)
        1
        (* x (fact (+ x -1)))))))
\end{ildcode}\end{listingbox}
\lstcaption{Example module \code{"demos/fact.ild"}}\label{example-module-fact}

The module loader's operating mechanism is as follows:
\begin{enumerate}
\item Given a source filename containing a module, call \ild{read-source} to obtain its
content as a value
\item \paraphrase{Verify that the expression in the source file is actually a module}
\item Obtain the list of imports and list of definitions \paraphrase{via standard implementations of keyword list operations}
\item Recursively process imported modules and build a binding environment \ild{env}
that contains all imports
\item Build a list of name-body pairs from the definitions listed in the module
\item Convert the module to a \ild{letrec}-expression \paraphrase{containing}
the module's definitions and using them to generate
list of key-value pairs for the module exports.
\end{enumerate}

Running the module loader on \cref{example-module-fact} would result in a program
equivalent to the following letrec expression:
% MTODO: letrec expression


\begin{note}
the following will be removed soon
\end{note}

In \cref{module-loader-outer} we construct a program called
\ild{module-loader} that is able to process
user programs organised in modules such as the one defined in
\cref{example-module-fact}.
\begin{listingbox}\begin{ildcode}
(!(eval (free-vars) (read-source
    "core/bootstrap/let.ild"))
    ( (host-bindings (free-vars))
      (letrec (eval (free-vars) (read-source
        "core/bootstrap/letrec.ild"))))
  (!letrec
      ( <definitions of helpers> )
    (!lambda (mod-name func-name args)
      (apply
        (kv-get func-name
          (load-module mod-name))
        args))))
\end{ildcode}\end{listingbox}
\lstcaption{Outermost structure of the module loader}\label{module-loader-outer}
The outermost \ild{let}-expression makes the root environment
(captured via \ildsf{free-vars}) and \ild{letrec} available to the
internal \ild{letrec} expression.
The module loader, when called with a module name,
entrypoint function name and a list of arguments, loads that module, picks the entrypoint function
from its export list and applies it to the given arguments.
The core helper is a function called \ild{eval-module},
whose implementation we show in \cref{def-eval-module} while
omitting technical details such as standard helpers
(list and \ild{kv-*} operations, \ild{concat},
\ild{lambda}, \ild{map}, \ild{if}, \ild{let*}).
The \emph{eval-module} function:
\begin{enumerate}
\item Recursively processes imported modules and builds a binding environment \ild{env}
that contains all imports
\item Builds a list of name-body pairs from the definitions listed in the module
\item Calls \ild{letrec} (as an ordinary function) on these definition pairs and
a body that generates a list of key-value pairs for the module exports.
\end{enumerate}
\begin{listingbox}\begin{ildcode}
; ---- inside letrec block ----
(load-module (!lambda (module-name)
  (eval-module (read-source module-name))))

(eval-module (!lambda (mod-expr)
  (!let* (
      (mod (cdr mod-expr))

      (imports (kv-fetch 'imports (list) mod))
      (defs    (kv-fetch 'defs    (list) mod))
      (exports (kv-fetch 'exports (list) mod))

      (env (process-imports imports))
      (eval-def (!lambda (def)
        (eval env def)))
      (def-items (map eval-def defs))

      (body (cons list
        (map
          (!lambda (sym)
            (list cons (list quote sym) sym))
          exports))))

    (eval env (letrec def-items body)))))
\end{ildcode}\end{listingbox}
\lstcaption{Definitions of \ild{load-module} and \ild{eval-module}}\label{def-eval-module}
The \ild{process-imports} helper (\cref{def-process-imports}) calls \ild{load-module} recursively on
imported modules and then merges the resulting export lists. For the special
case when the \ild{host} keyword is used in place of a module name,
the host bindings (obtained via the outermost call to \ild{free-vars}) are
used in place of an export list.
We omit the implementation of \ild{process-import-items}, which merely
picks the requested key-value pairs from the export lists.
\begin{listingbox}\begin{ildcode}
; ---- inside letrec block ----
(process-imports (!lambda (import-list)
  (concat (map process-import import-list))))

(process-import (!lambda (import-spec)
  (!if (!and (symbol? (car import-spec))
             (sym-eq? (car import-spec) 'host))
    (process-import-items
      host-bindings
      (cdr import-spec))
    (process-import-items
      (load-module (car import-spec))
      (cdr import-spec)))))
\end{ildcode}\end{listingbox}
\lstcaption{Definition of \ild{process-imports}}\label{def-process-imports}
Running a program
such as the one in \cref{example-module-fact} is now a matter of passing its
name and entrypoint to the module loader (shown in \cref{module-load-example}).
\begin{listingbox}\begin{ildcode}
((eval (free-vars) (read-source
  "core/bootstrap/module-loader.ild"))

  "demos/fact.ild"
  'main
  '())
\end{ildcode}\end{listingbox}
\lstcaption{Executing the module loader on a user program}\label{module-load-example}
The ILD interpreter and/or build system need not be concerned with any details
about loading and importing modules and can be kept simple: all of the logic
is implemented in ILD itself.
